% =====================================================================
%  Section 5: Discussion
% =====================================================================
\section{Discussion}
\label{sec:discussion}

\paragraph{What TADS delivers.}
\textsc{TADS} adds a representation-level structural signal to
dynamic instruction-tuning data selection. The no-anchor base ranks
candidates by the reward--advantage diagnostic consensus score
$\mathcal{R}_i^{(t)}a_i^{(t)}$; \textsc{TADS} keeps this base fixed
and multiplies it by a trajectory-anchor factor. Thus
$\lambda=0$ recovers the same consensus base, while $\lambda>0$
isolates the contribution of hidden-state trajectory alignment.

The key contribution is the anchor, not a new selection objective. The
anchor is model-intrinsic, requiring no seed examples, validation
labels, external judges, or auxiliary scorers. It is also
stability-controlled: under a positive eigengap, its movement is
bounded at the learning-rate scale. The guarantee therefore concerns
the structural component introduced by \textsc{TADS}, not the
reward--advantage product itself.

\paragraph{Why the anchor helps.}
Dynamic selection based only on scalar reward or advantage signals can
be unstable in multi-task instruction pools. In particular, the
variance-ratio weight used to combine difficulty and uncertainty can
drift with task composition because it mixes within-task and
between-task variance (Eq.~\ref{eq:total-variance}). The trajectory
anchor adds a pool-level geometric signal: it favours samples aligned
with the representation direction the model is currently shaping.

This also separates \textsc{TADS} from dynamic data selection for
RLVR, such as DOTS~\cite{sun2025dots}. DOTS selects GRPO questions
using rollout outcomes and verifiable rewards; \textsc{TADS}
selects supervised instruction--response pairs using hidden-state
geometry. DOTS is therefore an important related method, but not a
directly runnable baseline for our SFT setting, where automatic
verifiers, binary rewards, and rollout groups are unavailable by
default.

\paragraph{Implications and limitations.}
The stability bound makes the anchor a reproducible structural
object rather than probe noise. The eigengap $\gamma^{(t)}$ also
serves as a diagnostic: when it is large, the anchor is sharply
identified; when the covariance becomes nearly isotropic, the anchor
weakens and the method naturally reduces toward the no-anchor
consensus base.

\textsc{TADS} currently targets text-only SFT with single-sequence
instruction--response pairs and a rank-one anchor per layer. Future
extensions include turn-aware anchors for multi-turn dialogue,
separate anchors for multimodal representations, and rank-$r$
eigenspace anchors for samples combining multiple skills. Extending
trajectory anchoring to RL fine-tuning would require a different
experimental setup, since prompts, sampled responses, rollouts, and
reward-conditioned states define different possible anchor spaces.
